\chapter{Creating Professional Tables}
\label{chap:tables}
\index{table}\index{booktabs package}\index{tabularx package}\index{decimal alignment}

A table is not a storage container for every recorded value. It is an argument
in rows and columns: the reader should understand what is being compared, which
unit applies, and what each summary means. We begin with the smallest grid and
add only the structures required for a scholarly table.

\begin{learningoutcomes}
  \item Build a \code{tabular} with rows, columns, and alignment.
  \item Place a table, caption, and label in a float.
  \item Use \pkg{booktabs} rules and avoid unnecessary vertical lines.
  \item Create flexible and multi-line columns with \pkg{tabularx}.
  \item Align numerical columns with \pkg{siunitx}.
  \item Prepare table notes and a multi-page \code{longtable}.
  \item Explain a safe optional route from CSV data to a checked table.
\end{learningoutcomes}

\section{The smallest tabular}

The \code{tabular} environment aligns cells. Its required argument describes
the columns: \code{l}, \code{c}, and \code{r} mean left, centred, and right.

\begin{latexcode}{Complete example: a two-column tabular}
\documentclass{article}
\begin{document}
\begin{tabular}{lr}
Method & Complete pairs \\
A & 12 \\
B & 12
\end{tabular}
\end{document}
\end{latexcode}

Ampersands separate cells; \code{\\} ends a row. Every row needs the same
column logic. The environment aligns content but supplies neither a caption nor
a number.

\begin{commonerror}
\textbf{Symptom:} ``Extra alignment tab has been changed to \cmd{cr}.''

\textbf{Cause:} a row contains too many ampersands for the declared columns.

\textbf{Repair:} for two columns, confirm that each row has one ampersand.
\end{commonerror}

\section{Rules that organise rather than cage}

Load \pkg{booktabs} and use \cmd{toprule}, \cmd{midrule}, and \cmd{bottomrule}:

\begin{table}[htbp]
  \centering
  \caption{Descriptive summary of simulated paired measurements.}
  \label{tab:descriptive-summary}
  \begin{tabular}{@{}lSS@{}}
    \toprule
    {Method} & {Mean (mm)} & {SD (mm)} \\
    \midrule
    A & 5.68 & 2.71 \\
    B & 5.74 & 2.70 \\
    \bottomrule
  \end{tabular}
\end{table}

The \code{@\{\}} at each edge removes extra outer padding. The three rules
show the table boundary and distinguish the header; they do not draw a box
around every cell. Vertical rules and a rule after every row usually interrupt
the horizontal reading path.

\section{The table float, caption, and label}

The \code{table} environment is a \term{float}: an object \LaTeX{} may move to a
good page position. The optional \optional{htbp} offers here, top, bottom, or a
float page; it is a preference list, not a command to violate page design.

Place \cmd{caption} before \cmd{label} so the label records the table number.
Use a descriptive name such as \code{tab:descriptive-summary}. Refer to the
table from the prose so the reader knows why it matters.

\begin{goodpractice}
A caption should identify the population or material, variables, units, and
important qualifications. Do not use ``Results'' as a caption. Keep detailed
interpretation in the prose rather than hiding it beneath the table.
\end{goodpractice}

\section{Flexible text columns}

The \pkg{tabularx} environment takes a width and lets \code{X} columns share
the remaining space:

\begin{latexcode}{Complete example: a descriptive table}
\documentclass{article}
\usepackage{booktabs,tabularx}
\begin{document}
\begin{tabularx}{\textwidth}{@{}lX@{}}
\toprule
Item & Description \\
\midrule
Identifier & Stable code assigned to one simulated pair. \\
Difference & Method B minus Method A, reported in millimetres. \\
\bottomrule
\end{tabularx}
\end{document}
\end{latexcode}

For a fixed-width paragraph column, use \code{p\{width\}}. It wraps text and
top-aligns the cell. Choose the width from the page and content, then inspect
the result; a table with many paragraph columns may be better redesigned.

\section{A statistical table}

\begin{table}[htbp]
  \centering
  \caption{Illustrative analysis of simulated differences. Confidence limits
  are included only to demonstrate structure.}
  \label{tab:statistical-summary}
  \begin{tabular}{@{}lS[table-format=1.2]S[table-format=-1.2]S[table-format=1.2]@{}}
    \toprule
    {Quantity} & {Estimate} & {Lower limit} & {Upper limit} \\
    \midrule
    Mean difference (mm) & 0.06 & -0.18 & 0.30 \\
    SD of differences (mm) & 0.31 & {} & {} \\
    \bottomrule
  \end{tabular}

  \begin{minipage}{0.86\textwidth}
    \smallskip\footnotesize
    \textit{Note.} Values are simulated. The interval is not evidence about a
    real instrument or population.
  \end{minipage}
\end{table}

Short table notes can sit in a minipage aligned with the table. A journal class
may supply a dedicated notes environment; prefer it when available. An empty
braced cell protects an \code{S} column from non-numeric content.

\section{Tables across pages}

Use \code{longtable} for a table that must break across pages. It is not placed
inside a \code{table} float.

\begin{latexcode}{Longtable pattern with repeated headings}
\begin{longtable}{@{}lSS@{}}
\caption{Complete simulated measurements.}\label{tab:all-pairs}\\
\toprule
{Pair} & {Method A (mm)} & {Method B (mm)} \\
\midrule
\endfirsthead
\multicolumn{3}{l}{\tablename\ \thetable\ continued}\\
\toprule
{Pair} & {Method A (mm)} & {Method B (mm)} \\
\midrule
\endhead
\midrule
\multicolumn{3}{r}{Continued on next page}\\
\endfoot
\bottomrule
\endlastfoot
P01 & 1.20 & 1.40 \\
P02 & 2.10 & 1.90 \\
% Continue with verified or explicitly simulated rows.
\end{longtable}
\end{latexcode}

The first head appears once; the later head repeats. Foot definitions control
continuation text and the final rule. Test with enough rows to force a page
break, and confirm that the caption, number, header, and notes remain clear.

\section{Optional: data from CSV}

Copying a large dataset into source invites transcription errors. Packages can
read CSV files, or a reproducible analysis script can export a finished
\file{.tex} table. This is an advanced workflow because imported values still
need validation, rounding, escaping, captions, and version control.

\begin{advancedtopic}
Keep raw data separate from presentation. Record the script and software
version that created the table. Inspect the generated source and PDF. Never
assume that successful import proves that columns have the intended units or
that missing values were handled correctly.
\end{advancedtopic}

\begin{scienceinpractice}
A geological sample table may need sample code, location, lithology, depth,
and analytical result. If descriptions dominate, use flexible text columns; if
measurements dominate, use decimal-aligned columns. If the full record is very
long, show a concise analytical table in the chapter and place the complete
record in an appendix or repository.
\end{scienceinpractice}

\section{Code lab: turn observations into a table}

\begin{codelab}
Compile the complete file, replace one simulated value, and add a third method.
Then remove the caption's word ``Simulated'' and explain why that would be
scientifically misleading even though the source still compiles.
\end{codelab}

\begin{latexcode}{A compact, publication-style numerical table}
\documentclass{article}
\usepackage{booktabs}
\usepackage{siunitx}

\begin{document}
\begin{table}[htbp]
  \centering
  \caption{Simulated recovery measurements.}
  \label{tab:recovery}
  \begin{tabular}{l S[table-format=2.1] S[table-format=1.2]}
    \toprule
    Method & {Mean (\%)} & {SD} \\
    \midrule
    A & 91.4 & 2.10 \\
    B & 94.8 & 1.73 \\
    \bottomrule
  \end{tabular}
\end{table}
\end{document}
\end{latexcode}

\chapterproject

\begin{miniproject}
Replace the temporary numerical check with a numbered table containing the
simulated Method A mean, Method B mean, and mean difference. Use
\pkg{booktabs}, \code{S} columns, one shared unit in each heading, and a note
stating that all values are simulated. Label and mention the table in Results.
\end{miniproject}

\chaptersummary

\code{tabular} aligns cells; a \code{table} float adds placement, caption, and
numbering. \pkg{booktabs} uses a small hierarchy of horizontal rules.
\pkg{tabularx} and paragraph columns manage text; \pkg{siunitx} aligns numbers.
Notes define qualifications without replacing results prose. \code{longtable}
supports repeated headings across pages. Imported tables remain subject to
validation and visual inspection.

\chaptercheck
\begin{chapterchecklist}
  \item Each row has the declared number of columns.
  \item The caption says what the table contains and which data are simulated.
  \item The label follows the caption and has a \code{tab:} prefix.
  \item Units, rounding, missing values, and notes are unambiguous.
  \item Numerical columns align at the decimal marker.
  \item I did not use vertical rules merely by habit.
\end{chapterchecklist}

\chapterexercisestitle
\begin{chapterexercises}
  \item Build a two-column table with three data rows.
  \item Diagnose a row with one ampersand too many.
  \item Add \pkg{booktabs} rules and remove vertical rules.
  \item Convert a long description column to \code{tabularx}.
  \item Align positive and negative values with an \code{S} column.
  \item Write a caption and note for explicitly simulated environmental data.
  \item Explain why \code{longtable} is not placed inside \code{table}.
  \item Challenge: redesign a real table from your field for one clear reader
        task and document every omitted column.
\end{chapterexercises}
